% corpus_supplement: Erdős problem #551
% source: https://www.erdosproblems.com/latex/551
% fetched_at: 2026-07-14T18:47:52Z
% payload_sha256: 4eaf7b5907c0e85b442e1e170200fe4278565183f966d293803fbab22a98e2e4
% statement/remarks/references extracted by erdos_ingest.extract_source_page;
% rendered by erdos_ingest.render_tex — do not edit
\documentclass{article}
\usepackage{amsmath, amssymb, amsthm}
\usepackage{hyperref}
\usepackage{geometry}
\geometry{margin=1in}

\title{Erd\H{o}s Problem \#551}
\date{Last updated: 2025-08-31}
\author{Source: erdosproblems.com}

\begin{document}
\maketitle

\noindent\textbf{Status:} OPEN \quad \textbf{No prize}

\noindent\textbf{Tags:} graph theory, ramsey theory

\medskip
\noindent\textbf{Problem Statement:}

Prove that\[R(C_k,K_n)=(k-1)(n-1)+1\]for $k\geq n\geq 3$ (except when $n=k=3$).

\bigskip
\noindent\textbf{Remarks:}

Asked by Erd\H{o}s, Faudree, Rousseau, and Schelp, who also ask for the smallest value of $k$ such that this identity holds (for fixed $n$). They also ask, for fixed $n$, what is the minimum value of $R(C_k,K_n)$?

This identity was proved for $k>n^2-2$ by Bondy and Erd\H{o}s \cite{BoEr73}. Nikiforov \cite{Ni05} extended this to $k\geq 4n+2$.

Keevash, Long, and Skokan \cite{KLS21} have proved this identity when $k\geq C\frac{\log n}{\log\log n}$ for some constant $C$, thus establishing the conjecture for sufficiently large $n$.

This problem is #18 in Ramsey Theory in the graphs problem collection.

\bigskip
\noindent\textbf{References:}

[BoEr73] Bondy, J. A. and Erd\H{o}s, P., Ramsey numbers for cycles in graphs. J. Combinatorial Theory Ser. B (1973), 46-54.
 
 [KLS21] Keevash, Peter and Long, Eoin and Skokan, Jozef, Cycle-complete Ramsey numbers. Int. Math. Res. Not. IMRN (2021), 277-302.
 
 [Ni05] Nikiforov, V., The cycle-complete graph Ramsey numbers. Combin. Probab. Comput. (2005), 349-370.

\bigskip
\noindent\small{Source: \url{https://www.erdosproblems.com/551}}
\end{document}
